
\documentclass[11pt]{article}
\usepackage[margin=1in]{geometry}
\usepackage{amsmath, amssymb, amsthm}
\usepackage{hyperref}
\usepackage{enumitem}

\title{The Parasite Machine (Updated):\\
Asymmetry, Self-Referential Logic, and Conic Regimes}
\author{}
\date{}

\newtheorem{definition}{Definition}

\begin{document}
\maketitle

\section*{Abstract}
This document presents an expanded and refined formalization of the Parasite machine.
We integrate:
\begin{itemize}
\item asymmetric relations as locally generated intensities,
\item implied counter-relations via representational regimes,
\item the tetralemma as a self-referential dynamical system,
\item and conic geometry as curvature of implication.
\end{itemize}

\section{1. Ontological Commitment}

We begin with a refusal.

We do not assume:
\begin{itemize}
\item substances,
\item conserved quantities,
\item global symmetry,
\item global asymmetry.
\end{itemize}

Instead we assume only:
\begin{itemize}
\item positions,
\item relations,
\item transformations,
\item observables.
\end{itemize}

\paragraph{Narrative}
The system does not describe a world. It generates one.
Structure is not given—it emerges from dynamics.

\section{2. The Pre-Relational Field}

Before any directed relation, we posit:

\[
\mathcal{R}_0(A,B)
\]

\paragraph{Narrative}
This is not a relation in the usual sense.
It is an underdetermined field of possible relation.

It contains no direction, no symmetry, and no asymmetry.

\section{3. Generation of Asymmetry}

\begin{definition}
\[
R(A,B) = \Psi_C(\mathcal{R}_0(A,B))
\]
\end{definition}

\paragraph{Narrative}
Asymmetry is not primitive.

It is generated locally by the action of a third term $C$.
The parasite produces direction by acting on an undetermined field.

Thus:
\[
A \xrightarrow[\iota,\delta]{} B
\]
is a stabilization, not a foundation.

\section{4. Implied Counter-Relation}

\begin{definition}
\[
\Phi_m : R(A,B) \mapsto R(B,A)
\]
\end{definition}

\paragraph{Narrative}
Reciprocity is not given.
It is implied, and the form of implication depends on a mode $m$.

\section{5. Representation as Curvature}

\[
\mathcal{M} = \{Id, Res, An, Con\}
\]

\[
\kappa(m) \in \{ellipse, parabola, hyperbola, mixed\}
\]

\paragraph{Narrative}
Representation is not correspondence.
It is curvature.

Each mode determines how a relation folds back onto itself.

\begin{itemize}
\item Identity closes (ellipse)
\item Resemblance drifts (parabola)
\item Analogy transforms (mixed)
\item Contrariety diverges (hyperbola)
\end{itemize}

\section{6. The Parasite}

\begin{definition}
\[
C \triangleright_m R(A,B)
\]
\end{definition}

\paragraph{Narrative}
The parasite acts on:
\begin{itemize}
\item intensity,
\item asymmetry,
\item implied relation,
\item representational regime,
\item logical state.
\end{itemize}

It is not external.
It is the condition under which relations take form.

\section{7. The Tetralemma (Dynamical)}

\[
\mathcal{T} = \{affirmed,\ negated,\ both,\ neither\}
\]

\paragraph{Narrative}
The tetralemma is not a classification of truth.

It is a state space.

\subsection*{Crucial Shift}
We do not select a corner.

Instead:

\[
\tau_{t+1} = \tau(\tau_t, R, \Phi, C)
\]

\paragraph{Narrative}
The system traverses the tetralemma.

Symmetry and asymmetry are not global properties.
They are regimes that appear and disappear through this traversal.

\section{8. Conic Regimes and Logic}

\paragraph{Narrative}
The tetralemma and conics are not separate.

They describe the same structure from different perspectives.

\begin{itemize}
\item Ellipse: closure, mutual implication, stability of ``both''
\item Parabola: threshold, drift, one-sided implication
\item Hyperbola: divergence, separation, instability
\item Mixed: transformation, mediation, ambiguity
\end{itemize}

Thus:
logic = geometry = dynamics

\section{9. Full Machine}

\begin{definition}
\[
P = (A,B,C,R,\Phi,m,\tau,d)
\]
\end{definition}

\[
P' = \Psi(P)
\]

\paragraph{Narrative}
Everything evolves.

At each step:
\begin{itemize}
\item a relation is generated,
\item a counter-relation is implied,
\item a tetralemmatic state is proposed,
\item that state is transformed,
\item positions are updated,
\item roles may permute,
\item regimes may shift.
\end{itemize}

\section{10. Computational Instantiation}

The Python implementation realizes this machine explicitly.

\paragraph{Narrative}
The code demonstrates:
\begin{itemize}
\item asymmetry as dynamic,
\item reciprocity as derived,
\item logic as evolving,
\item representation as regime,
\item and the parasite as operator.
\end{itemize}

The distinction between:
\begin{itemize}
\item candidate tetralemma state,
\item updated tetralemma state
\end{itemize}
is critical.

It shows that logic is not read off the system—it is transformed by it.

\section{11. Conceptual Synthesis}

\begin{itemize}
\item There is no global symmetry.
\item There is no global asymmetry.
\item The system does not resolve this.
\item It continuously rearticulates it.
\end{itemize}

\paragraph{Final Statement}

A world emerges from:
\begin{itemize}
\item a pre-relational field,
\item local generation of asymmetry,
\item mode-dependent implication,
\item self-referential traversal of logical regimes,
\item and parasitic mediation.
\end{itemize}

\end{document}
